\paragraph{分辨率塔的 Haar 极限与 Tate 模块中的秩分裂}
在跨分辨率仿射刚性已经锁定之后，H.159 的 dyadic 分辨率塔与 H.160 的规范 Gödel 半直积仿射表示在同一 Tate 模块中还强制出一条更细的测度--秩--几何分裂链：有限扭点层的均匀分布在逆极限上闭合为 full-rank Haar 几何，而任何规范语义地址像仍只能坍缩到由单一地址向量生成的一条 rank-one \(B\)-进 Cantor 子系统。下列诸条把这一分裂压缩为可直接调用的闭式结论。

\begin{theorem}[分辨率塔的 Haar 极限律]\label{thm:xi-time-part9gc-resolution-tower-haar-limit}
设
\[
T:=T_2(E_{\min})
\cong
\varprojlim_n E_{\min}[2^n]
\cong
\ZZ_2^2
\]
为推论 \ref{cor:killo-2adic-tate-hypercube-limit} 给出的逆极限识别。对每个 \(n\ge 1\)，记 \(\mu_n\) 为有限群 \(E_{\min}[2^n]\) 上的均匀概率测度。则：
\begin{enumerate}
  \item 在乘二投影
  \[
  [2]:E_{\min}[2^{n+1}]\longrightarrow E_{\min}[2^n]
  \]
  下，测度族 \((\mu_n)_{n\ge 1}\) 严格相容；
  \item 存在唯一 Borel 概率测度 \(\mu_\infty\) 于 \(T\) 上，使得对每个 \(n\) 与每个模 \(2^n\) 余类
  \[
  C(a,b;n):=
  \left\{
  x\in T:\ x\equiv (a,b)\pmod{2^nT}
  \right\},
  \qquad
  (a,b)\in (\ZZ/2^n\ZZ)^2,
  \]
  都有
  \[
  \mu_\infty\bigl(C(a,b;n)\bigr)=4^{-n};
  \]
  \item \(\mu_\infty\) 正是紧群 \(T\cong \ZZ_2^2\) 的归一化 Haar 概率测度。
\end{enumerate}
\leanverified{paper\_xi\_time\_part9gc\_resolution\_tower\_haar\_limit}
\end{theorem}
\begin{proof}
由推论 \ref{cor:killo-2adic-tate-hypercube-limit}，每一层都可坐标化为
\[
E_{\min}[2^n]\cong (\ZZ/2^n\ZZ)^2,
\]
且乘二映射恰对应模约化
\[
(a\bmod 2^{n+1},\,b\bmod 2^{n+1})
\longmapsto
(a\bmod 2^n,\,b\bmod 2^n).
\]
因此，对每个深度 \(n\) 的余类 \(C(a,b;n)\)，其在下一层恰有 \(4\) 个 lift；于是
\[
\mu_{n+1}\bigl([2]^{-1}\{(a,b)\}\bigr)
=
4\cdot 4^{-(n+1)}
=
4^{-n}
=
\mu_n\bigl(\{(a,b)\}\bigr),
\]
这就证明了相容性。

现以 cylinder 质量
\[
\mu_\infty\bigl(C(a,b;n)\bigr):=4^{-n}
\]
定义逆极限上的柱测度。该定义是相容的，因为每个深度 \(n\) 的 cylinder 正好分裂为 \(4\) 个深度 \(n+1\) 的 cylinder。故由 cylinder 扩张唯一得到 \(T\) 上的 Borel 概率测度 \(\mu_\infty\)。

最后，由 \(T\cong \ZZ_2^2\) 可知每个 \(C(a,b;n)\) 都是开子群 \(2^nT\) 的一个陪集，而这类陪集在任意平移下彼此置换。由于 \(\mu_\infty\) 在同一深度的全部 \(4^n\) 个陪集上取相同质量 \(4^{-n}\)，故 \(\mu_\infty\) 在全部基本 clopen 陪集上都满足平移不变性。此类集合生成 \(T\) 的 Borel \(\sigma\)-代数，故 \(\mu_\infty\) 即为归一化 Haar 概率测度。证毕。
\end{proof}

\begin{corollary}[分辨率塔的 \texorpdfstring{$2$}{2}-进秩二刚性]\label{cor:xi-time-part9gc-resolution-tower-rank-two-rigidity}
\leanverified{paper\_xi\_time\_part9gc\_resolution\_tower\_rank\_two\_rigidity}
设 \(M\) 为自由 \(\ZZ_2\)-模，且存在对一切 \(n\ge 1\) 都与模 \(2^n\) 商相容的层同构
\[
M/2^nM\cong E_{\min}[2^n].
\]
若 \(M\cong \ZZ_2^r\)，则必有
\[
r=2.
\]
特别地，任何试图以单一 \(2\)-进地址轴 \(\ZZ_2\) 忠实实现整条分辨率塔的方案都不可能成立。
\end{corollary}
\begin{proof}
若 \(M\cong \ZZ_2^r\)，则
\[
\#(M/2^nM)=2^{rn}.
\]
另一方面，由推论 \ref{cor:killo-2adic-tate-hypercube-limit}，
\[
E_{\min}[2^n]\cong (\ZZ/2^n\ZZ)^2,
\]
故
\[
\#E_{\min}[2^n]=2^{2n}.
\]
逐层比较得 \(2^{rn}=2^{2n}\) 对全部 \(n\) 成立，只能有 \(r=2\)。若 \(r=1\)，则 \(\#(M/2^nM)=2^n\) 与 \(\#E_{\min}[2^n]=4^n\) 矛盾，故单轴实现不可能。证毕。
\end{proof}

\begin{theorem}[规范 Gödel 编码的 \texorpdfstring{$B$}{B}-进 Cantor 吸引子与前缀轨道闭包]\label{thm:xi-time-part9gc-canonical-godel-cantor-attractor-orbit-closure}
沿用定理 \ref{thm:killo-godel-semidir-affine-2adic} 与定理
\ref{thm:xi-time-part9g-canonical-submonoid-free-monoid} 的记号。固定 primitive 地址向量
\[
v\in T\setminus 2T,
\qquad
B:=2^L>M,
\qquad
T_v:=\ZZ_2v\subset T.
\]
定义
\[
\Xi_{B,M}:\{0,\dots,M\}^{\NN}\longrightarrow T_v,
\qquad
\Xi_{B,M}(a_1,a_2,\dots):=\sum_{t\ge 1}a_tB^{t-1}v,
\]
以及其像
\[
K_{B,M}(v):=\Xi_{B,M}\bigl(\{0,\dots,M\}^{\NN}\bigr).
\]
则成立：
\begin{enumerate}
  \item \(\Xi_{B,M}\) 是到 \(K_{B,M}(v)\) 的拓扑同胚；等价地，每个 \(x\in K_{B,M}(v)\) 都具有唯一的 \(B\)-进 digit 展开；
  \item \(K_{B,M}(v)\) 满足严格不交并
  \[
  K_{B,M}(v)=\bigsqcup_{a=0}^{M}\bigl(av+B K_{B,M}(v)\bigr);
  \]
  \item 若 \(H_{\mathrm{can}}\) 表示规范编码子单子，则
  \[
  \overline{\{\Psi(r,k)(0):\ (r,k)\in H_{\mathrm{can}}\}}
  =
  K_{B,M}(v).
  \]
\end{enumerate}
\end{theorem}
\begin{proof}
取有限事件集
\[
E:=\{0,\dots,M\},
\qquad
\mathrm{code}:E\hookrightarrow \NN,\quad \mathrm{code}(a):=a.
\]
则由定理 \ref{thm:xi-time-part62d-hologram-image-cantor-homeomorphism}，地址映射 \(X:E^{\NN}\to \Lambda_{E,v}\) 在此正好就是 \(\Xi_{B,M}\)，且其像满足严格不交并
\[
\Lambda_{E,v}
=
\bigsqcup_{a=0}^{M}\bigl(av+B\Lambda_{E,v}\bigr).
\]
于是第 \(1\) 与第 \(2\) 条成立，并且 \(K_{B,M}(v)=\Lambda_{E,v}\)。

现证第 \(3\) 条。任取
\[
x=\Xi_{B,M}(a_1,a_2,\dots)\in K_{B,M}(v).
\]
对每个 \(k\ge 1\)，定义 \(r_k\in \mathcal R\) 为
\[
r_k(p_t):=
\begin{cases}
a_t,& 1\le t\le k,\\
0,& t>k.
\end{cases}
\]
则 \((r_k,k)\in H_{\mathrm{can}}\)，并且
\[
\Psi(r_k,k)(0)=\sum_{t=1}^{k}a_tB^{t-1}v.
\]
由于 \(|B|_2<1\)，上述有限前缀和在 \(k\to\infty\) 时收敛到 \(x\)。故
\[
K_{B,M}(v)\subseteq
\overline{\{\Psi(r,k)(0):\ (r,k)\in H_{\mathrm{can}}\}}.
\]
反向包含是显然的，因为每个 \(\Psi(r,k)(0)\) 都是某条无限 digit 序列在尾部补零后得到的地址点，故本身属于 \(K_{B,M}(v)\)。证毕。
\end{proof}

\begin{theorem}[规范 Gödel 吸引子的精确维数与 Haar 测度二歧律]\label{thm:xi-time-part9gc-canonical-godel-attractor-dimension-haar-dichotomy}
沿用定理 \ref{thm:xi-time-part9gc-canonical-godel-cantor-attractor-orbit-closure} 的记号。把 \(T_v\cong \ZZ_2\) 视作一维 \(2\)-进紧群，则
\[
\dim_{\mathrm H}K_{B,M}(v)
=
\dim_{\mathrm M}K_{B,M}(v)
=
\frac{\log(M+1)}{\log B}
=
\frac{\log_2(M+1)}{L}.
\]
此外：
\begin{enumerate}
  \item 若 \(M=B-1\)，则
  \[
  K_{B,M}(v)=T_v;
  \]
  \item 若 \(M\le B-2\)，则
  \[
  \mu_{T_v}\bigl(K_{B,M}(v)\bigr)=0,
  \qquad
  \operatorname{Int}_{T_v}\bigl(K_{B,M}(v)\bigr)=\varnothing,
  \]
  其中 \(\mu_{T_v}\) 为 \(T_v\) 的归一化 Haar 概率测度。
\end{enumerate}
\end{theorem}
\begin{proof}
由定理 \ref{thm:xi-time-part9gc-canonical-godel-cantor-attractor-orbit-closure}，集合 \(K_{B,M}(v)\) 与
\[
\Lambda_{E,v}
\qquad
\bigl(E=\{0,\dots,M\},\ \mathrm{code}=\mathrm{id}\bigr)
\]
完全一致。因而，结论 \ref{con:xi-time-part62b-rankone-hologram-selfsimilar-measure-dimension} 的第 \(3\) 条给出维数公式，而结论 \ref{con:xi-time-part62c-godel-tate-minimal-bitwidth-threshold} 的第 \(2\)、\(3\) 条给出满地址线与真 Cantor 子集的二歧律。于是
\[
\dim_{\mathrm H}K_{B,M}(v)
=
\dim_{\mathrm M}K_{B,M}(v)
=
\frac{\log|E|}{\log B}
=
\frac{\log(M+1)}{\log B}.
\]
若 \(M=B-1\)，则结论 \ref{con:xi-time-part62c-godel-tate-minimal-bitwidth-threshold} 的第 \(2\) 条给出 \(K_{B,M}(v)=T_v\)。若 \(M\le B-2\)，则其第 \(3\) 条给出 Haar 零测；再由结论 \ref{con:xi-time-part62b-rankone-hologram-selfsimilar-measure-dimension} 的第 \(2\) 条，内部必为空。证毕。
\end{proof}

\begin{theorem}[规范仿射字值的忠实阈值恰为 \texorpdfstring{$C>M$}{C>M}]\label{thm:xi-time-part9gc-canonical-affine-word-faithful-threshold}
\leanverified{paper\_xi\_time\_part9gc\_canonical\_affine\_word\_faithful\_threshold}
设 \(C=2^s\ge 2\)，并定义有限字集
\[
\mathcal W_M:=\bigsqcup_{k\ge 0}\{0,\dots,M\}^k.
\]
对每个字
\[
\mathbf a=(a_1,\dots,a_k)\in \{0,\dots,M\}^k
\]
定义仿射字值
\[
\Xi_C(\mathbf a)(x):=
C^k x+\sum_{t=1}^{k}a_tC^{t-1}v
\qquad (x\in T_v).
\]
则映射
\[
\Xi_C:\mathcal W_M\longrightarrow \operatorname{Aff}(T_v)
\]
忠实当且仅当
\[
C>M.
\]
\end{theorem}
\begin{proof}
先设 \(C>M\)。若
\[
\Xi_C(a_1,\dots,a_k)=\Xi_C(b_1,\dots,b_\ell)
\]
为同一仿射映射，则分别在 \(0\) 与 \(v\) 处取值并相减，得到
\[
C^k v=C^\ell v.
\]
由于 \(T_v\) 对 \(\ZZ_2\) 无挠且 \(v\neq 0\)，必有 \(C^k=C^\ell\)，从而 \(k=\ell\)。再由
\[
\sum_{t=1}^{k}(a_t-b_t)C^{t-1}v=0
\]
可知
\[
\sum_{t=1}^{k}(a_t-b_t)C^{t-1}=0
\qquad\text{于 }\ZZ_2.
\]
若存在首个分歧位置 \(j\) 使 \(a_j\neq b_j\)，则上式可写为
\[
C^{j-1}\bigl((a_j-b_j)+Cu\bigr)=0
\]
对某个 \(u\in \ZZ_2\) 成立。由于
\[
0<|a_j-b_j|\le M<C,
\]
有 \((a_j-b_j)\notin C\ZZ_2\)，故 \((a_j-b_j)+Cu\neq 0\)，矛盾。于是所有 digits 都相同，故 \(\Xi_C\) 忠实。

反之，若 \(C\le M\)，则长度 \(2\) 的两字
\[
(C,0),
\qquad
(0,1)
\]
都属于 \(\{0,\dots,M\}^2\)，并且对一切 \(x\in T_v\) 都有
\[
\Xi_C(C,0)(x)=C^2x+C v=\Xi_C(0,1)(x).
\]
故 \(\Xi_C\) 不忠实。证毕。
\end{proof}

\begin{theorem}[规范 Gödel 仿射层的精确字增长指数]\label{thm:xi-time-part9gc-canonical-affine-semigroup-word-growth}
\leanverified{paper\_xi\_time\_part9gc\_canonical\_affine\_semigroup\_word\_growth}
沿用前述记号。对每个 \(k\ge 0\)，记
\[
H_{\mathrm{can}}(k):=\{(r,k)\in H_{\mathrm{can}}\},
\qquad
\mathcal N_k:=\bigl|\{\Psi(r,k):\ (r,k)\in H_{\mathrm{can}}(k)\}\bigr|.
\]
则有精确公式
\[
\mathcal N_k=(M+1)^k.
\]
因而规范仿射层的字增长熵
\[
h_{\mathrm{grow}}
:=
\lim_{k\to\infty}\frac{1}{k}\log \mathcal N_k
\]
满足
\[
h_{\mathrm{grow}}=\log(M+1).
\]
特别地，结合定理 \ref{thm:xi-time-part9gc-canonical-godel-attractor-dimension-haar-dichotomy}，
\[
\dim_{\mathrm H}K_{B,M}(v)=\frac{h_{\mathrm{grow}}}{\log B}.
\]
\end{theorem}
\begin{proof}
由定理 \ref{thm:xi-time-part9g-canonical-submonoid-free-monoid}，每个 \(H_{\mathrm{can}}(k)\) 的元素都唯一对应一个长度 \(k\) 的 digit 字
\[
(a_1,\dots,a_k)\in \{0,\dots,M\}^k.
\]
而在 \(C=B\) 的情形下，定理 \ref{thm:xi-time-part9gc-canonical-affine-word-faithful-threshold} 表明该长度 \(k\) 的字值映射是忠实的。故不同 digit 字给出不同仿射映射，遂有
\[
\mathcal N_k=\bigl|\{0,\dots,M\}^k\bigr|=(M+1)^k.
\]
两边取对数并除以 \(k\)，即得
\[
h_{\mathrm{grow}}=\log(M+1).
\]
最后一式只是把定理 \ref{thm:xi-time-part9gc-canonical-godel-attractor-dimension-haar-dichotomy} 的维数公式重写为
\[
\frac{\log(M+1)}{\log B}
=
\frac{h_{\mathrm{grow}}}{\log B}.
\qedhere
\]
\end{proof}

\begin{theorem}[同一 Tate 环境中的 Haar--Cantor 分裂]\label{thm:xi-time-part9gc-haar-cantor-splitting}
记
\[
T:=T_2(E_{\min})\cong \ZZ_2^2,
\qquad
\mu_{\mathrm{res}}:=\mu_\infty,
\]
其中 \(\mu_\infty\) 为定理 \ref{thm:xi-time-part9gc-resolution-tower-haar-limit} 给出的 Haar 极限测度。对任意 \(M<B\)，设 \(\nu\) 为任意支撑在 \(K_{B,M}(v)\) 上的 Borel 概率测度。则
\[
\nu\perp \mu_{\mathrm{res}}.
\]
\end{theorem}
\begin{proof}
由定理 \ref{thm:xi-time-part9gc-canonical-godel-cantor-attractor-orbit-closure}，
\[
K_{B,M}(v)\subseteq T_v=\ZZ_2v.
\]
因此只需证明
\[
\mu_{\mathrm{res}}(T_v)=0.
\]
对每个 \(n\ge 1\)，记自然约化映射为
\[
\rho_n:T\longrightarrow T/2^nT\cong (\ZZ/2^n\ZZ)^2.
\]
由于 \(T_v\) 是由单一向量 \(v\) 生成的循环 \(\ZZ_2\)-子模，其像 \(\rho_n(T_v)\) 在有限群 \(T/2^nT\) 中仍是循环子群，因而
\[
\#\rho_n(T_v)\le 2^n.
\]
另一方面，由定理 \ref{thm:xi-time-part9gc-resolution-tower-haar-limit}，每个模 \(2^nT\) 陪集的 \(\mu_{\mathrm{res}}\)-质量都等于 \(4^{-n}\)。故
\[
\mu_{\mathrm{res}}(T_v)
\le
\#\rho_n(T_v)\cdot 4^{-n}
\le
2^n\cdot 4^{-n}
=
2^{-n}.
\]
令 \(n\to\infty\) 得 \(\mu_{\mathrm{res}}(T_v)=0\)。

现在取可测集 \(A:=T_v\)。则
\[
\nu(A)=1,
\qquad
\mu_{\mathrm{res}}(A)=0.
\]
按互相奇异的定义，遂得 \(\nu\perp \mu_{\mathrm{res}}\)。证毕。
\end{proof}

\begin{corollary}[语义 Cantor 子系统对 full-rank 采样空间的仿射不可恢复性]\label{cor:xi-time-part9gc-affine-nonrecoverability-fullrank-space}
沿用定理 \ref{thm:xi-time-part9gc-haar-cantor-splitting} 的记号。对任意连续仿射自映射
\[
A:T\longrightarrow T,
\qquad
A(x)=Lx+t,
\]
都有
\[
A\bigl(K_{B,M}(v)\bigr)\subseteq t+L(T_v).
\]
其中 \(L(T_v)\) 是秩至多为 \(1\) 的闭 \(\ZZ_2\)-子模。因而
\[
A\bigl(K_{B,M}(v)\bigr)
\]
在 \(T\) 中始终 Haar 零测且内部为空，特别不可能等于 \(T\)。
\end{corollary}
\begin{proof}
由定理 \ref{thm:xi-time-part9gc-canonical-godel-cantor-attractor-orbit-closure}，
\[
K_{B,M}(v)\subseteq T_v.
\]
故
\[
A\bigl(K_{B,M}(v)\bigr)\subseteq A(T_v)=t+L(T_v).
\]
而 \(T_v\) 是循环 \(\ZZ_2\)-模，其在线性映射 \(L\) 下的像 \(L(T_v)\) 仍是循环 \(\ZZ_2\)-模；又因 \(L\) 连续且 \(T_v\) 紧，\(L(T_v)\) 闭。故 \(L(T_v)\) 的秩至多为 \(1\)。

若 \(L(T_v)=\{0\}\)，则 \(A(K_{B,M}(v))\) 为单点，结论显然。若 \(L(T_v)\neq \{0\}\)，则与定理 \ref{thm:xi-time-part9gc-haar-cantor-splitting} 的证明完全相同，可得任意 rank-one 闭子模在 \(T\cong \ZZ_2^2\) 中都是 Haar 零测。于是其仿射平移 \(t+L(T_v)\) 亦为 Haar 零测。

在紧群 \(T\) 中，任一非空开集都具有正 Haar 测度；因此 Haar 零测集不可能含有内点。故
\[
A\bigl(K_{B,M}(v)\bigr)
\subseteq
t+L(T_v)
\]
也必内部为空，特别不可能等于整个 \(T\)。证毕。
\end{proof}

\begin{theorem}[binary Gödel 吸引子上的 \texorpdfstring{$\varphi$}{phi}-logistic 维数曲线]\label{thm:xi-time-part9gc-binary-godel-logistic-dimension-curve}
沿用定理 \ref{thm:xi-time-part9gc-canonical-godel-cantor-attractor-orbit-closure} 的记号，并取二元字母情形 \(M=1\)。记
\[
\Xi_B:=\Xi_{B,1}:\{0,1\}^{\NN}\longrightarrow K_{B,1}(v).
\]
对每个 \(q\ge 0\)，定义
\[
\theta(q):=\frac{1}{1+\varphi^{q+1}},
\]
并令 \(\beta_q\) 为 \(\{0,1\}^{\NN}\) 上的 Bernoulli 乘积测度，
\[
\beta_q(\omega_1=1)=\theta(q),
\qquad
\beta_q(\omega_1=0)=1-\theta(q).
\]
再记
\[
\nu_q:=(\Xi_B)_*\beta_q,
\qquad
h(t):=-t\log t-(1-t)\log(1-t).
\]
则 \(\nu_q\) 为 exact-dimensional 测度，并且
\[
\dim_{\mathrm H}\nu_q=\frac{h(\theta(q))}{\log B}.
\]
此外，对任意 \(p,q\ge 0\)，escort 族的极限 KL 散度满足
\[
\lim_{m\to\infty}D_{\mathrm{KL}}(\pi_{m,q}\Vert \pi_{m,p})
=
D_{\mathrm{KL}}\!\bigl(\mathrm{Bernoulli}(\theta(q))\Vert \mathrm{Bernoulli}(\theta(p))\bigr).
\]
\end{theorem}
\begin{proof}
在 \(E=\{0,1\}\)、\(\mathrm{code}=\mathrm{id}\) 的情形，定理
\ref{thm:xi-time-part9gc-canonical-godel-cantor-attractor-orbit-closure} 表明
\[
K_{B,1}(v)=\Xi_B(\{0,1\}^{\NN}),
\]
且 \(\Xi_B\) 为拓扑同胚。

对每个 \(q\)，把
\[
\mathbf p_q:=(1-\theta(q),\,\theta(q))
\]
视为 \(\{0,1\}\) 上的概率向量。则 \(\beta_q\) 正是定理 \ref{thm:xi-time-part62dc-hologram-bernoulli-exact-dimensional} 中的 Bernoulli 乘积测度 \(\nu_{\mathbf p_q}\)，而 \(\nu_q\) 正是其地址推前 \(\mu_{\mathbf p_q}\)。故同一定理立即给出 \(\nu_q\) 的 exact-dimensional 性以及
\[
\dim_{\mathrm H}\nu_q
=
\frac{H(\mathbf p_q)}{\log B}.
\]
由于
\[
H(\mathbf p_q)
=
-\theta(q)\log\theta(q)
-(1-\theta(q))\log(1-\theta(q))
=
h(\theta(q)),
\]
遂得
\[
\dim_{\mathrm H}\nu_q=\frac{h(\theta(q))}{\log B}.
\]

另一方面，定理 \ref{thm:xi-time-part9s-escort-logistic-expfamily-bregman} 已把
\[
q\longmapsto \theta(q)=\frac{1}{1+\varphi^{q+1}}
\]
识别为 escort 极限族的 logistic 指数坐标，而定理 \ref{thm:xi-time-part9od-escort-escort-fdiv-binary-closure} 在 \(f(u)=u\log u\) 时恰给出
\[
\lim_{m\to\infty}D_{\mathrm{KL}}(\pi_{m,q}\Vert \pi_{m,p})
=
\bigl(1-\theta(q)\bigr)\log\frac{1-\theta(q)}{1-\theta(p)}
+\theta(q)\log\frac{\theta(q)}{\theta(p)}.
\]
右端正是
\[
D_{\mathrm{KL}}\!\bigl(\mathrm{Bernoulli}(\theta(q))\Vert \mathrm{Bernoulli}(\theta(p))\bigr).
\]
故第二条成立。证毕。
\end{proof}

\noindent
定理 \ref{thm:xi-time-part9gc-resolution-tower-haar-limit} 至定理 \ref{thm:xi-time-part9gc-binary-godel-logistic-dimension-curve} 共同表明，H.159 的分辨率塔极限并不与 H.160 的规范语义地址像混同：前者在 \(T_2(E_{\min})\cong \ZZ_2^2\) 上携带 full-rank 的 Haar 几何，后者则始终困在由单一地址向量生成的 rank-one \(B\)-进 Cantor 子系统内。于是，同一 Tate 环境中的分辨率采样协议与 Gödel 语义协议不是简单的坐标改写关系，而是两类在秩、测度与仿射可恢复性上同时分裂的读出机制。

\endinput
